\documentclass[a4paper,11pt]{article}
\usepackage[utf8]{inputenc}
\usepackage[T1]{fontenc}
\usepackage[french]{babel}
\usepackage[left=2cm,right=2cm,top=2cm,bottom=2cm]{geometry}
\usepackage{xcolor}
\usepackage{hyperref}
\usepackage{fontawesome5}
\usepackage{parskip}

% --- Couleurs Thales ---
\definecolor{primary}{RGB}{0, 45, 114}
\hypersetup{colorlinks=true, urlcolor=primary}

\begin{document}

\noindent
\begin{minipage}[t]{0.5\textwidth}
    \textbf{NJIHA KUITO Brayanne Stella}\\
    Elève Ingénieure -- Master EEA\\[3pt]
    \faMapMarker\ Cergy (95), France\\
    \faPhone\ (+33) 6 52 31 08 46\\
    \faEnvelope\ \href{mailto:stellanjiha1@gmail.com}{stellanjiha1@gmail.com}
\end{minipage}
\hfill
\begin{minipage}[t]{0.4\textwidth}
    \raggedleft
    Cergy, le 1 juin 2026
\end{minipage}

\vspace{1.2cm}

\hfill
\begin{minipage}[t]{0.5\textwidth}
    \raggedleft
    \textbf{Thales -- Bureau d'Etudes Circuits Imprimés}\\
    Service Recrutement\\
    Elancourt, France
\end{minipage}

\vspace{1cm}

\noindent\textbf{Objet :} Candidature au poste d'Alternante -- Technicienne Projeteur Cartes Electroniques -- Alternance 1 an

\vspace{0.8cm}

Madame, Monsieur,

La réalisation de schémas électroniques, l'implantation des composants et le routage de circuits imprimés sont des activités que je pratique depuis ma première année à l'ENSEA. Contribuer aux projets d'antennes pour systèmes aéroportés et spatiaux au sein du Bureau d'Etudes Circuits Imprimés de Thales à Elancourt représente pour moi une opportunité concrète d'approfondir ces compétences dans un environnement d'excellence technologique. C'est pourquoi je vous adresse ma candidature pour cette alternance à partir de septembre 2026, dans le cadre de mon Master 2 EEA à CY Paris Université.

Dans le cadre d'un projet de robot à navigation autonome à l'ENSEA, j'ai assuré la conception complète d'un PCB sous KiCad : saisie du schéma de principe, implantation des composants et routage d'une carte d'alimentation en intégrant les contraintes CEM, thermiques et de puissance. J'ai ensuite constitué le dossier de définition complet (fichiers XAO, plans, nomenclature) avant de procéder à l'assemblage, la soudure de précision et la mise au point du prototype en laboratoire, en étroite collaboration avec les équipes de conception pour garantir la cohérence des choix techniques. Cette expérience de bout en bout m'a donné une vision claire des différentes étapes du développement d'une carte électronique, de la définition jusqu'à la validation.

Mes travaux de caractérisation CEM, portant sur l'analyse des couplages électromagnétiques et l'impact des plans de masse sur l'immunité des signaux, m'ont appris à travailler avec méthode dans un contexte de validation exigeant, en interaction avec les contraintes liées aux technologies hyperfréquences que Thales développe. Ma pratique de la simulation sous LTSpice et PSpice complète cette approche en me permettant d'anticiper les comportements des circuits avant prototypage.

Je suis disponible dès septembre 2026 pour une alternance d'un an et serais ravie de rejoindre l'équipe du Bureau d'Etudes Circuits Imprimés de Thales Elancourt.

Je reste à votre disposition pour un entretien afin de discuter de ma candidature.

Je vous prie d'agréer, Madame, Monsieur, l'expression de mes salutations distinguées.

\vspace{0.8cm}

\textbf{NJIHA KUITO Brayanne Stella}\\
\textit{Elève Ingénieure ENSEA / CY Paris Université}

\end{document}
